\section{Conclusion}
\subsection{Summary of Findings}

\subsubsection{Empirical Confirmation of the Barren Plateau}
\textbf{[PENDING -- to be finalized once the AI4I n=8 benchmark is available; see Section~6.1.]} This project designed, implemented, and evaluated AL-QNN, a variational quantum classifier whose depth is scheduled online by a reinforcement-learning agent, validated at up to 8 qubits on the hardware available to the team (Section~1.6). An earlier draft of this section reported a barren-plateau confirmation on the qubit axis and a primary AI4I benchmark measured at 12 qubits; both have been withdrawn because they fall outside the project's 8-qubit scope and, independently, because the AI4I data loader has since been corrected (Section~6.1). They are pending re-measurement within the 8-qubit range.

\subsubsection{Scheduler Performance Across Regimes}
\textbf{[PENDING -- see note above; this section's strongest supporting evidence, a 20-qubit local-cost regime, has been removed as out of scope.]} Under regimes where gradients are already healthy (\Cref{tab:heart,tab:sim8g}) a static heuristic -- greedy -- matches or beats the learned agents, with PPO tracking a shallow fixed-depth baseline rather than showing a clear advantage. The one regime in this report where the learned scheduler outright wins is the supplementary German Credit deployment (\Cref{tab:german}, single-seed, Section~6.1). The learned scheduler is therefore not shown to be universally superior at the qubit counts evaluated here; whether it separates from static baselines more clearly at $n=8$ on the primary AI4I benchmark is the open question this section will answer once \Cref{tab:ai4i} is re-run.

\subsection{Contributions and Reflections on the Project}

\subsubsection{Contributions}
The project contributes: (i) a telemetry-driven MDP formulation of quantum-circuit depth scheduling with a sixteen-dimensional state and a five-action space; (ii) a practical two-simulator architecture (surrogate for training, PennyLane for deployment) that makes RL training tractable for quantum circuits; (iii) an evenhanded, multi-regime evaluation that identifies the conditions in which adaptive scheduling is worthwhile; and (iv) a reproducible, log-backed methodology.

\subsubsection{Reflections on the Process}
Reflecting on the process, the two-simulator design and the single-source-of-truth ansatz were the decisions that most improved both speed and correctness, while the main difficulty was resisting overclaiming from unstable single-seed runs.

\subsection{Threats to Validity}
Four factors limit how far the conclusions in Section~6 generalize, and are stated here explicitly rather than left implicit.

\subsubsection{Single-Seed Variance (Internal Validity)}
The primary AI4I benchmark (\Cref{tab:ai4i}) is the only three-seed table in Section~6, pending re-run at $n=8$; every other scheduler-comparison table in Section~6 is a single run. The extended fixed-depth sweep (Section~6.2) shows at least one clearly non-monotonic curve (the synthetic q = 8, global-cost regime), which is direct evidence that single-seed numbers can and do contain noise large enough to change which depth looks best; the same could in principle be true of any single headline accuracy figure elsewhere in the report.

\subsubsection{No Held-Out Test Split (Construct Validity)}
The deployment and comparison scripts behind every table in Section~6 split each dataset into a 70/30 train/validation partition only, with no separate held-out test set (Section~4.5); every reported accuracy, F1, and ROC AUC figure is therefore a validation-split metric, computed on the same data the scheduler observes every epoch to make its accept/rollback decisions. This is not a data-leakage bug -- no scaler, PCA, or model parameter is ever fit on this split -- but it does mean the reported numbers should be read as in-loop validation performance rather than performance on data the pipeline never touched, and absolute accuracy figures may carry a mild optimistic bias relative to a true held-out test set.

\subsubsection{Surrogate Fidelity (Construct Validity)}
Every ``simulate''-labelled result is generated by a classical surrogate calibrated on a limited set of real measurements (Section~4.7.1); its trends are reliable within the calibrated regime but are not a guarantee of behaviour on the real circuit outside it, which is precisely why the PennyLane-real \Cref{tab:heart} is retained as an independent anchor rather than being replaced by a cheaper synthetic equivalent.

\subsubsection{Regime Coverage (External Validity)}
All reported regimes sit at a single qubit count ($n=8$), the ceiling validated on available hardware (Section~1.6); within that fixed qubit count, two cost types (local and global) are compared, supplemented by three single-seed real deployments on other datasets (German Credit, Parkinson's, Pima; Section~6.1) reported for completeness rather than as additional matched regimes. This isolates the cost-type factor but says nothing about how either scheduler's advantage would change at a higher qubit count, since that axis was cut from the study entirely rather than sampled; the recommendations in Section~6.4 should be read as guidance for where to look next, not as a closed characterization of the full regime space, and any claim about qubit-count scaling beyond $n=8$ is outside what this report can support.

\subsection{Limitations and Future Work}
The principal limitations are that results are obtained on a noiseless simulator rather than physical hardware, that evaluation was capped at 8 qubits by the RAM and runtime cost of PennyLane statevector simulation on available hardware (Section~1.6), and that several secondary results are reported with a single seed. The most immediate item of future work is completing the AI4I benchmark at $n=8$ (\Cref{tab:ai4i}) and the corresponding barren-plateau scaling measurement, both currently pending (Section~6.1); beyond that, the most important item is to complete multi-seed runs for every regime so that the cross-regime claims can be stated with confidence intervals, and, hardware permitting, to extend the qubit range past 8 to test whether the cost-type effect observed here strengthens at larger qubit counts.

